\section{异常}

\begin{question}
	以下哪个异常类用于处理数组下标越界异常？
	\begin{options}
		\item ArrayIndexOutOfBoundsException
		\item NullPointerException
		\item ArithmeticException
		\item ClassCastException
	\end{options}
	\begin{answer}
		A
	\end{answer}
	\begin{explanation}
		\texttt{ArrayIndexOutOfBoundsException} 表示访问数组时索引超出范围。
	\end{explanation}
\end{question}

\begin{question}
	在Java中，以下哪个关键字用于捕获异常？
	\begin{options}
		\item try
		\item catch
		\item throw
		\item throws
	\end{options}
	\begin{answer}
		B
	\end{answer}
	\begin{explanation}
		\texttt{catch} 用于捕获并处理异常。
	\end{explanation}
\end{question}

\begin{question}
	以下关于Java中异常处理中finally块的说法正确的是
	\begin{options}
		\item 无论是否发生异常都会执行
		\item 只有发生异常才会执行
		\item 只有不发生异常才会执行
		\item 只有在try块中有return语句时才会执行
	\end{options}
	\begin{answer}
		A
	\end{answer}
	\begin{explanation}
		\texttt{finally} 块总会被执行（除非 JVM 退出）。
	\end{explanation}
\end{question}

\begin{question}
	以下哪种情况会导致OutOfMemoryError？
	\begin{options}
		\item 数组越界
		\item 除数为0
		\item 内存泄漏
		\item 空指针异常
	\end{options}
	\begin{answer}
		C
	\end{answer}
	\begin{explanation}
		内存泄漏导致内存耗尽，最终抛出 \texttt{OutOfMemoryError}。
	\end{explanation}
\end{question}
